\section{Porovnání možností nasazení}
\label{sec:selection-deployment}

Jazykové modely lze do systému integrovat pomocí dvou hlavních přístupu, a to buď jako externí službu poskytovanou třetí stranou, nebo jako lokálně hostovaný model běžící přímo na infrastruktuře školy.
Tato sekce tak analyzuje obě varianty z pohledu technické realizace, provozních nákladů, bezpečnosti dat a dlouhodobé udržitelnosti.

\subsection{Cloudové API}
\label{subsec:deployment-cloud}

Cloudové API je způsob přístupu k jazykovému modelu provozovanému na serverech externího poskytovatele.
Integrace probíhá prostřednictvím HTTP požadavků, systém odešle dotaz na API endpoint a následně obdrží odpověď s vygenerovanou odpovědí.
Výpočet tak probíhá úplně na straně poskytovatele, což znamená, že není potřeba žádná specializovaná hardware pro provoz modelu a škálování je zajištěno poskytovatelem služby.

Mezi nejznámější postykovatele cloudových API pro LLM patří:
\begin{itemize}
    \item \textbf{OpenAI} -- poskytuje přístup k modelům GPT-5.4, GPT-5-mini a mnohé další varianty\cite{openai_api}.
    \item \textbf{Anthropic} -- nabízí jako hlavní modely Claude Opus~4.6, Claude Sonnet~4.6 a Claude Haiku~4.5\cite{anthropic_api}.
    \item \textbf{Google} -- přístup k modelům Geminy prostřednictvím Google AI Studio nebo Vertex AI~\cite{google_api}.
    \item \textbf{Mistral AI} -- poskytuje přístup k vlastním modelům Mistral Large a Mistral Nemo, které jsou k dispozici i jako open-source ke stažení.~\cite{mistral_api}.
\end{itemize}

\subsubsection*{Hardwarové nároky}

Z pohledu provozování systému kelvin jsou hardwarové nároky pro použití cloudového API prakticky nulové.
Veškerý výpočet probíhá na serverech poskytovatele a na straně systému kelvin je potřeba pouze stabilní internetové připojení a server, který zvládne tyto HTTP požadavky zpracovat.
Jelikož školní infrastruktura leží na optické síti CETIN, není tak téměř žádný problém zajistit dostatečnou konektivitu pro komunikaci s cloudovým API\@.

\subsubsection*{Limity}

Cloudové API zavádějí několik klíčových typů omezení, které je nutné při návrhu systému zohlednit:

\begin{enumerate}
    \item \textbf{Limit na počet tokenů v kontextu} -- každý model má maximální délku zpracovatelného vstupu i výstupu.
    Toto omezení je zpravidla u nejmodernějších modelů (state-of-the-art models) dostatečné, běžně 128~000 tokenů nebo více, ale je důležité zajistit, aby systém Kelvin dokázal limitovat požadavky tak, aby nebyl limit překročen -- překročení limitu API signalizuje chybovou odpovědí, která musí být na straně systému správně zachycena a ošetřena.

    \item \textbf{Limit na počet požadavků} -- poskytovatelé často zavádějí takzvané rate limits, které omezují počet požadavků, které lze poslat v určitém časovém období (například 100 požadavků za minutu).
    Pro školní prostředí, kde může být v určitých časech (například před termínem odevzdání) zvýšený počet požadavků, je důležité tyto limity zohlednit a případně implementovat mechanismy pro řízení zátěže (rate limiting) na straně systému Kelvin\@.

    \item \textbf{Finanční náklady} -- cloudové API jsou zpoplatněny na základě počtu zpracovaných tokenů, což může být pro rozsáhlé použití nákladné.
    Je důležité pečlivě odhadnout očekávané náklady na základě počtu studentů, frekvence odevzdání a průměrné velikosti odevzdání, aby bylo možné zajistit udržitelnost provozu.

    \item \textbf{Závislost na třetí straně} -- použití cloudového API znamená, že systém Kelvin bude závislý na dostupnosti a spolehlivosti služby poskytovatele.
    V případě výpadku služby nebo změny v API může být funkčnost systému narušena, což je riziko, které je třeba zvážit.
\end{enumerate}

\subsubsection*{Ochrana dat}

Použití cloudového API se záměrem k ochraně dat studentů bylo rozebráno v sekci~\ref{subsubsec:gdpr-data-transfer}.
Tato ochrana je často tak klíčová, že může být rozhodujícím faktorem pro volbu mezi cloudovým API a on-premise nasazením.

\subsubsection*{Náklady}

Náklady cloudových API jsou určeny cenovým modelem \textit{za token}.
Poskytovatelé zpravidla rozlišují vstupní tokeny (text odesílaný modelu) a výstupní tokeny (text generovaný modelem), přičemž výstupní tokeny jsou typicky dražší.

\begin{table}[H]
    \centering
    \resizebox{\textwidth}{!}{%
        \begin{tabular}{llll}
            \toprule
            \textbf{Model}      & \textbf{Vstupní tokeny (\$/1M)} & \textbf{Výstupní tokeny (\$/1M)} & \textbf{Přibližný náklad} \\
            \midrule
            GPT-5.4 (OpenAI)    & 2,50                            & 15,00                            & $\sim$\TODO{Naklady} \$   \\
            GPT-5-mini (OpenAI) & 0,25                            & 2,00                             & $\sim$\TODO{Naklady} \$   \\
            GPT-5-nano (OpenAI) & 0,05                            & 0,40                             & $\sim$\TODO{Naklady} \$   \\
            Claude 4.6 Opus     & 5,00                            & 25,00                            & $\sim$\TODO{Naklady} \$   \\
            Claude 4.6 Sonnet   & 3,00                            & 15,00                            & $\sim$\TODO{Naklady} \$   \\
            Claude 3.5 Haiku    & 0,80                            & 4,00                             & $\sim$\TODO{Naklady} \$   \\
            Claude 3.0 Haiku    & 0,25                            & 1,25                             & $\sim$\TODO{Naklady} \$   \\
            Gemini 3.1 Pro      & 2,00/4,00                       & 5,00                             & $\sim$\TODO{Naklady} \$   \\
            Gemini 3.1 Flash    & 0,50                            & 3,00                             & $\sim$\TODO{Naklady} \$   \\
            Gemini 3 Flash      & 12,00/18,00                     & 0,30                             & $\sim$\TODO{Naklady} \$   \\
            \bottomrule
        \end{tabular}
    }
    \caption{Přibližné ceny vybraných cloudových API modelů \cite{openai_api, anthropic_api, google_vertex} pro odevzdání s $\sim$5~000 vstupními tokeny a $\sim$1~000 výstupními tokeny}
    \label{tab:cloud-api-costs}
\end{table}

\TODO{Doplnit přibližný počet vstupních a výstupních tokenu pro jedno odevzdání -> vypočítat náklady}
% TODO: Footnote pro Geminy, první číslo je prompts <= 200k tokens, druhé číslo je prompts > 200k tokens
% TODO: Footnote pro OpenAI, liší pro typ licence, Standard je dražší než Flex

\TODO{Závěr, který z těch modelů se jeví jako nejvhodnější z pohledu nákladů a výkonu pro školní prostředí, případně zdůvodnit, že je potřeba provést další testování pro konkrétní odhad nákladů.}

\subsection{On-premise řešení}
\label{subsec:deployment-on-premise}

On-premise nasazení znamená provozování jazykového modelu přímo na lokálních serverech, bez využití externích cloudových služeb.
V tomto případě je model hostován na serverové infrastruktuře a veškeré zpracování dat probíhá lokálně.

Výhodou on-premise řešení je především nezávislost na externích poskytovatelích služeb a absence poplatků za jednotlivé API požadavky.
Správci mají zároveň možnost plně kontrolovat konfiguraci modelu, způsob jeho aktualizace i způsob škálování výpočetních prostředků.
Na druhou stranu je nutné zajistit dostatečné hardwarové prostředky, zejména výkonné grafické akcelerátory (GPU), dostatek operační paměti a vhodnou infrastrukturu pro provoz a správu modelů.

Pro usnadnění provozu jazykových modelů v lokálním prostředí existuje několik softwarových nástrojů, které poskytují rozhraní pro jejich správu, spouštění a inferenci:

\begin{itemize}
    \item \textbf{Ollama}~\cite{ollama} -- nástroj zaměřený na jednoduché spouštění a správu jazykových modelů v lokálním prostředí.
    Poskytuje jednotné rozhraní pro práci s různými modely a umožňuje jejich snadné stažení, spuštění i aktualizaci.
    Součástí nástroje je také API, které umožňuje integraci modelů do externích aplikací.

    \item \textbf{Groq}~\cite{groq} -- platforma zaměřená na vysoce výkonnou inferenci jazykových modelů.
    Nabízí optimalizovanou infrastrukturu pro rychlé zpracování požadavků nad velkými modely.
    Přestože je Groq často využíván jako cloudová služba, jeho technologie je založena na specializovaném hardwaru navrženém pro efektivní provoz modelů.

    \item \textbf{vLLM}~\cite{vllm} -- knihovna určená pro efektivní provoz velkých jazykových modelů, která se zaměřuje zejména na optimalizaci paměťového využití a paralelního zpracování požadavků.
    Využívá techniky jako \textit{PagedAttention}, které umožňují efektivně sdílet paměť mezi více požadavky a tím výrazně zvyšují propustnost systému.

    \item \textbf{llama.cpp}~\cite{llamacpp} -- open-source implementace zaměřená na provoz modelů rodiny LLaMA a jejich variant i na běžných procesorech bez nutnosti využití GPU\@.
    Nástroj umožňuje provoz kvantizovaných verzí modelů, což výrazně snižuje paměťové nároky a umožňuje nasazení i na méně výkonném hardwaru.
\end{itemize}

Volba konkrétního nástroje závisí především na dostupném hardwaru, požadované rychlosti inferenčního procesu a způsobu integrace do existující infrastruktury systému Kelvin.
Zatímco nástroje jako \textit{vLLM} jsou vhodné pro provoz ve výkonném serverovém prostředí s GPU akcelerací, řešení jako \textit{llama.cpp} umožňují provoz modelů i na méně výkonném hardwaru.
Nástroj \textit{Ollama} pak poskytuje kompromis mezi jednoduchostí nasazení a možností integrace do aplikací prostřednictvím API\@.

\subsubsection*{Hardwarové nároky}

Pro on-premise nasazení je zcela zásadní dostupnost vhodného hardware.
Inference jazykových modelů je výpočetně náročná operace (viz sekce~\ref{subsec:llm-computational-cost}), proto je využití GPU s dostatečnou kapacitou paměti VRAM téměř nezbytné pro přijatelné výkony.
přibližné nároky na hardware pro různé modely lze vidět v tabulce~\ref{tab:memory-requirements}.

\TODO{Doplnit porovnání CPU vs GPU pro inferenci. Spustit ollama na argexu a změřit výkon pro různé modely.}

\subsubsection*{Limity}

On-premise nasazení přináší vlastní sadu omezení:

\begin{itemize}
    \item \textbf{Omezení dostupným hardware} -- kvalita a rychlost analýzy jsou přímou funkcí dostupného GPU.
    Na nedostatečném hardware je nutné volit menší nebo silněji kvantizované modely, což může snižovat kvalitu výstupů.

    \item \textbf{Škálovatelnost} -- cloudové API automaticky škáluje s nárůstem zátěže, kdežto lokální server má fixní kapacitu.
    Při hromadném odevzdávání se požadavky řadí do fronty a čekají na zpracování.

    \item \textbf{Správa modelu} -- nové verze modelů musí být ručně staženy a nasazeny.
    Aktualizace nebo výměna modelu vyžaduje přímou intervenci správce systému.
\end{itemize}

\subsubsection*{Ochrana dat}

Z pohledu ochrany dat je on-premise nasazení výrazně jednodušší situace.
Studentské zdrojové kódy a zadání úloh nikdy neopustí infrastrukturu školy.
Neexistuje třetí strana, které by musela být svěřena správa dat, a není potřeba uzavírat DPA s externím poskytovatelem.
Z hlediska GDPR jde o přímočarou variantu: správce dat (škola) je totožný s provozovatelem infrastruktury.

\subsubsection*{Náklady}

Náklady on-premise nasazení mají odlišnou strukturu než u cloudového API:

\begin{itemize}
    \item \textbf{Jednorázové pořizovací náklady} -- GPU server vhodný pro inferenci (např.\ NVIDIA RTX 4090 s 24 GB VRAM) se pohybuje v ceně od přibližně 60~000 Kč výše.
    Pro výkonnější varianty (NVIDIA A100) jsou náklady řádově vyšší (200~000--500~000 Kč).

    \item \textbf{Průběžné provozní náklady} -- elektřina (moderní GPU má spotřebu 150--400 W při zatížení), chlazení serverovny a lidská práce nutná pro správu a aktualizaci systému.

    \item \textbf{Nulové náklady za token} -- po zaplacení hardware jsou marginalní náklady na každé zpracované odevzdání prakticky nulové, bez ohledu na počet odevzdání.
\end{itemize}

Pro systém s malým počtem uživatelů a odevzdání může být on-premise pořizovací investice ekonomicky nevýhodná ve srovnání s cloudovým API.
Naopak pro systém se stovkami studentů zpracovávajícími tisíce odevzdání ročně se on-premise investice může vyplatit v horizontu 1--2 let.

\begin{table}[H]
    \centering
    \resizebox{\textwidth}{!}{%
        \begin{tabular}{lll}
            \toprule
            \textbf{Vlastnost}   & \textbf{Cloudové API}              & \textbf{On-premise nasazení} \\
            \midrule
            Hardwarové nároky    & Žádné (výpočet u poskytovatele)    & GPU server (10+ GB VRAM)     \\
            Pořizovací náklady   & Nulové                             & Vysoké (GPU hardware)        \\
            Průběžné náklady     & Za token (opakující se)            & Elektřina a správa           \\
            Škálovatelnost       & Automatická (rate limity)          & Fixní kapacita               \\
            Ochrana dat          & Data opouštějí školu               & Data zůstávají na škole      \\
            Kvalita modelů       & Nejnovější (GPT-4o, Claude Sonnet) & Závisí na hardware           \\
            Latence              & Závislá na internetu a zátěži API  & Závisí na hardware           \\
            Správa a aktualizace & Automatická (poskytovatel)         & Manuální (správce)           \\
            \bottomrule
        \end{tabular}
    }
    \caption{Srovnání cloudového API a on-premise nasazení jazykového modelu z hlediska klíčových vlastností}
    \label{tab:deployment-comparison}
\end{table}

\endinput




\endinput